\documentclass[14pt,a4paper,oneside]{extbook}
\usepackage[utf8]{vietnam}
\usepackage{amsmath,amssymb,mathrsfs,fancyhdr,enumerate,multirow,makecell}
\usepackage{tkz-euclide,tikz-3dplot,tikz,tkz-tab}
\usetikzlibrary{arrows.meta,calc,intersections,angles,quotes,backgrounds,shapes.geometric,patterns,through,shadows,positioning}
\usepackage{yhmath}
\usepackage[top=0.7cm, bottom=0.7cm, left=1.0cm, right=1.0cm, headheight=16pt, headsep=3mm, footskip=5mm]{geometry}
\usepackage[hidelinks,unicode]{hyperref}
\usepackage{qrcode}
\usepackage[many]{tcolorbox}
\usepackage{fontawesome}

\pagestyle{fancy}
\fancyhf{}
\fancyhead[L]{\footnotesize\bfseries\color{teal!85!black}\faBook~SỔ TAY KIẾN THỨC TOÁN 4}
\fancyhead[R]{\footnotesize\bfseries\color{blue!85!black}Lớp Toán Thầy Lĩnh -- loptoanthaylinh.vn}
\fancyfoot[C]{\footnotesize\bfseries\color{gray!80!black}-- Sổ tay kiến thức: Bài 73. Ôn tập chung cuối năm --}
\renewcommand{\headrulewidth}{0.4pt}

\begin{document}

% --- BANNER TIÊU ĐỀ BÀI HỌC ---
\noindent
\begin{tcolorbox}[
  enhanced,
  colback=teal!8!white,
  colframe=teal!80!black,
  arc=3mm, boxrule=1pt,
  left=3mm, right=3mm, top=2mm, bottom=2mm,
  halign=center
]
  {\Large\bfseries\color{teal!90!black}\faBookmark\ BÀI 73: ÔN TẬP CHUNG CUỐI NĂM}\\[1.5pt]
  {\footnotesize\itshape\color{black!75} Đại tổng kết 4 mạch kiến thức cốt lõi Toán 4 chuẩn bị nền tảng lên lớp 5 trên 1 trang A4}
\end{tcolorbox}

\vspace{-2.5mm}

\noindent
\begin{minipage}[t]{0.49\textwidth}\vspace{0pt}
\begin{tcolorbox}[
  enhanced,
  colback=blue!4!white,
  colframe=blue!75!black,
  arc=2mm, boxrule=0.8pt,
  left=2mm, right=2mm, top=1.8mm, bottom=1.8mm,
  title={\footnotesize\bfseries\color{white}\faCalculator\ 1. Mạch Số học \& Tính toán}
]
{\footnotesize
\begin{itemize}\setlength{\itemsep}{1pt}\setlength{\leftskip}{-2.5mm}
  \item \textbf{Số tự nhiên:} Đọc, viết đến lớp triệu; làm tròn số; 4 phép tính $(+, -, \times, :)$ có nhớ và có dư.
  \item \textbf{Phân số:}
  \begin{itemize}\setlength{\itemsep}{0.5pt}\setlength{\leftskip}{-2.5mm}
    \item Rút gọn về tối giản \& Quy đồng mẫu số.
    \item 4 Phép tính với phân số.
    \item Tìm phân số của một số: $A \times \dfrac{m}{n}$.
  \end{itemize}
\end{itemize}
}
\end{tcolorbox}
\end{minipage}
\hfill
\begin{minipage}[t]{0.49\textwidth}\vspace{0pt}
\begin{tcolorbox}[
  enhanced,
  colback=teal!4!white,
  colframe=teal!75!black,
  arc=2mm, boxrule=0.8pt,
  left=2mm, right=2mm, top=1.8mm, bottom=1.8mm,
  title={\footnotesize\bfseries\color{white}\faCube\ 2. Mạch Hình học \& Đo lường}
]
{\footnotesize
\begin{itemize}\setlength{\itemsep}{1pt}\setlength{\leftskip}{-2.5mm}
  \item \textbf{Hình học:}
  \begin{itemize}\setlength{\itemsep}{0.5pt}\setlength{\leftskip}{-2.5mm}
    \item Hai đường thẳng vuông góc \& song song.
    \item Hình bình hành: $S = a \times h$.
    \item Hình thoi: $S = \dfrac{m \times n}{2}$.
  \end{itemize}
  \item \textbf{Đo lường:} Khối lượng (tấn, tạ, yến, kg), Diện tích ($m^2, dm^2, cm^2, mm^2$), Thời gian (thế kỉ, giây).
\end{itemize}
}
\end{tcolorbox}
\end{minipage}

\vspace{-2.5mm}
% --- KHỐI 3: SƠ ĐỒ TƯ DUY TIKZ ---
\noindent
\begin{tcolorbox}[
  enhanced,
  colback=orange!4!white,
  colframe=orange!80!black,
  arc=2mm, boxrule=0.8pt,
  left=2mm, right=2mm, top=1.5mm, bottom=1.5mm,
  title={\footnotesize\bfseries\color{white}\faSitemap\ SƠ ĐỒ TƯ DUY: CÂY ĐẠI THỤ KIẾN THỨC TOÁN 4}
]
\centering
\begin{tikzpicture}[
  box/.style={draw, rounded corners=2pt, align=center, font=\scriptsize\bfseries, inner sep=3pt, minimum width=3.4cm, minimum height=0.9cm},
  arr/.style={-{Stealth[length=2mm]}, thick}
]
  \node[draw=purple!80!black, fill=purple!15, rounded corners=3pt, font=\small\bfseries\color{purple!90!black}, inner sep=4pt, align=center] (core) at (5, 0.7) {TOÀN BỘ\\TOÁN 4};

  \node[box, fill=blue!10, draw=blue!80!black] (k1) at (0.8, 1.4) {1. SỐ TỰ NHIÊN\\Lớp triệu \& Phép tính};
  \node[box, fill=teal!10, draw=teal!80!black] (k2) at (9.2, 1.4) {2. PHÂN SỐ\\Bản chất \& 4 phép tính};
  \node[box, fill=orange!12, draw=orange!80!black] (k3) at (0.8, 0.0) {3. HÌNH HỌC\\Bình hành, hình thoi};
  \node[box, fill=green!15, draw=green!70!black] (k4) at (9.2, 0.0) {4. ĐO LƯỜNG \& THỐNG KÊ\\Đổi đơn vị, Biểu đồ cột};

  \draw[arr, blue!80!black] (core) -- (k1);
  \draw[arr, teal!80!black] (core) -- (k2);
  \draw[arr, orange!80!black] (core) -- (k3);
  \draw[arr, green!70!black] (core) -- (k4);
\end{tikzpicture}
\end{tcolorbox}

\vspace{-2.5mm}
% --- KHỐI 4: VÍ DỤ & MẸO TRÁNH BẪY ---
\noindent
\begin{minipage}[t]{0.49\textwidth}\vspace{0pt}
\begin{tcolorbox}[
  enhanced,
  colback=green!4!white,
  colframe=green!60!black,
  arc=2mm, boxrule=0.8pt,
  left=2mm, right=2mm, top=1.8mm, bottom=1.8mm,
  title={\footnotesize\bfseries\color{white}\faCheckCircle\ Ví dụ thực hành chuẩn mực}
]
{\footnotesize
\textbf{Tổng hợp cuối năm:}\\
Một mảnh đất hình thoi có hai đường chéo dài $14\text{ m}$ và $10\text{ m}$. Người ta dùng $\dfrac{2}{5}$ diện tích mảnh đất để trồng hoa. Tính diện tích trồng hoa?\\
\textbf{Bài giải:}\\
Diện tích mảnh đất là: $(14 \times 10) : 2 = 70\text{ m}^2$.\\
Diện tích trồng hoa là: $70 \times \dfrac{2}{5} = \mathbf{28}\text{ m}^2$.\\
\textbf{Đáp số:} $28\text{ m}^2$.
}
\end{tcolorbox}
\end{minipage}
\hfill
\begin{minipage}[t]{0.49\textwidth}\vspace{0pt}
\begin{tcolorbox}[
  enhanced,
  colback=red!4!white,
  colframe=red!70!black,
  arc=2mm, boxrule=0.8pt,
  left=2mm, right=2mm, top=1.8mm, bottom=1.8mm,
  title={\footnotesize\bfseries\color{white}\faExclamationTriangle\ Mẹo tránh bẫy trắc nghiệm}
]
{\footnotesize
\begin{itemize}\setlength{\itemsep}{1.5pt}\setlength{\leftskip}{-2.5mm}
  \item \textbf{Cẩm nang làm bài thi cuối năm:}\\
  1. Luôn đọc kĩ yêu cầu đề (tính giá trị, rút gọn hay so sánh).\\
  2. Kết quả phân số bắt buộc phải \textbf{tối giản}.\\
  3. Đo lường phải cùng \textbf{đơn vị đo} trước khi tính.\\
  4. Toán lời văn phải có \textbf{lời giải, phép tính, đơn vị và đáp số} rõ ràng.
\end{itemize}
}
\end{tcolorbox}
\end{minipage}

\end{document}
